\documentclass{tietreport}

% Import images
\usepackage{graphicx}

% Bibliography configuration
\usepackage[
  date=year,
  style=alphabetic,
  backend=biber,
  sorting=nty,
  maxnames=5,
  minnames=3,
  maxalphanames=4,
  minalphanames=3,
  backref=true,
  doi=false,
  isbn=false,
  url=false,
  eprint=false
]{biblatex}
\addbibresource{references.bib}

%% ==================================================
%% PROJECT INFORMATION
%% ==================================================

% Course/Lab header
\TitlePageHeader{%
  UCS503P - Project%
}

% Main project title
\ProjectTitle{%
  MediBridge
}

% Subtitle or project description
\ProjectSubTitle{%
  Healthcare Access and Financial Assistance Platform
}

% Student information
\author{%
  \texttt{1024160137} Kunal Motwani \\
  \texttt{1024160058} Suryansh Pahuja%
}

% Group number, degree, year, program
\TitlePageSubText{%
  Group: MediBridge, BE Third Year, CSE%
}

% Faculty advisor/guide name
\AdvisorName{%
  Dr. Sandeep Kaur%
}

% Evaluation period and department info
\TitlePageFooterText{{%
    \large Mid-Semester Evaluation \\[0.2cm]
    \bfseries Computer Science and Engineering
    Department%
  } \\[0.15cm]
  Thapar Institute of Engineering and Technology,
  Patiala%
}

%% ==================================================

\begin{document}

% Generate title page
\maketitle

% Table of contents
\tableofcontents

% ==================================================
% CHAPTER 1
% ==================================================

\chapter{Introduction}

\section{Background and Context}

Healthcare information is often fragmented across hospitals, doctors,
medical documents, and financial service providers. Patients may have
difficulty discovering suitable healthcare facilities, maintaining a
continuous record of their medical history, understanding treatment
costs, and arranging financial assistance when treatment becomes
expensive.

MediBridge was proposed as a centralized, patient-focused platform that
combines healthcare discovery, medical information management, and
financial assistance into a unified web application. The original
proposal identifies the core engineering objective as improving
healthcare accessibility, reducing delays, and supporting better
decision-making by bringing these workflows together.

The proposed system follows an MVP-first development approach. The
initial focus was placed on patient registration, healthcare profiles,
hospital discovery, medical record management, and progressively
extending the platform towards doctor interaction, financial assistance,
and emergency access.

\subsection{Problem Statement}

Patients commonly encounter the following difficulties:

\begin{itemize}
  \item Healthcare information about hospitals, doctors, treatment
    options, and costs is distributed across multiple sources.
  \item Medical records may be maintained separately by different
    hospitals, making continuity of care difficult.
  \item Patients may have limited visibility into who has accessed their
    medical information.
  \item Treatment costs can create a financial barrier even when
    suitable healthcare services are available.
  \item In emergencies, critical information such as blood group,
    allergies, medications, and emergency contacts may not be immediately
    accessible.
\end{itemize}

The project therefore aims to provide a single platform through which a
patient can manage health information, discover hospitals and doctors,
interact with healthcare providers, and access financial and emergency
support.

\section{Project Objectives}

The major objectives of the project, derived from the original proposal
and refined according to the implementation completed by the
mid-semester stage, are:

\begin{enumerate}
  \item \textbf{Develop a secure healthcare identity and authentication
  system:} provide patient and doctor registration, login, role-based
  access, password hashing, token-based authentication, and automatic
  generation of a unique MediBridge Health ID for patients.

  \item \textbf{Develop centralized medical information management:}
  allow patients to maintain profiles, store medical records, and
  access structured prescriptions while allowing authorized doctors to
  interact with patient information.

  \item \textbf{Implement healthcare discovery and appointment
  workflows:} provide hospital search, hospital details, doctor
  information, filtering, and appointment-request functionality.

  \item \textbf{Provide privacy-aware doctor access:} implement
  patient-controlled access through a PIN-based mechanism, temporary
  access sessions, access logs, and session revocation.

  \item \textbf{Integrate financial assistance workflows:} implement
  loan application storage, rule-based eligibility evaluation, and an
  external ML-model API interface for loan prediction.

  \item \textbf{Provide an emergency health-access mechanism:} store a
  limited emergency profile and expose an emergency-facing route that
  can be associated with a patient's Health ID and NFC-based workflow.
\end{enumerate}

% ==================================================
% CHAPTER 2
% ==================================================

\chapter{Methodology and Design}

\section{System Architecture}

The implemented MediBridge system follows the three-tier architecture
proposed in the project proposal, consisting of a React frontend,
Node.js/Express backend, and MongoDB-based persistence layer. An
additional Python Flask service is used for the financial prediction
component.

\subsection{High-Level Design}

The system is organized into the following major components:

\begin{enumerate}
  \item \textbf{Frontend:} a React single-page application built with
    Vite. The frontend provides separate patient and doctor workflows,
    protected routes, role-based navigation, profile management,
    healthcare discovery, record management, appointments, and financial
    application interfaces.

  \item \textbf{Backend API:} a Node.js and Express.js REST API that
    provides authentication, profile management, hospital search,
    medical-record operations, prescriptions, doctor access, financial
    services, fundraisers, emergency profiles, and appointment
    operations.

  \item \textbf{Database:} MongoDB accessed through Mongoose. The
    database uses dedicated models for users, hospitals, doctors,
    medical records, prescriptions, access sessions, access logs,
    emergency profiles, loan applications, fundraisers, donations, and
    appointments.

  \item \textbf{ML Service:} a separate Python Flask application
    exposes a \texttt{/predict} endpoint for loan prediction. The main
    backend communicates with this service through HTTP.

  \item \textbf{Cloud file storage:} uploaded profile and medical
    documents are integrated with Cloudinary through the backend.
\end{enumerate}

The overall design therefore preserves the proposal's separation between
presentation, business logic, and persistence while adding a separate
service boundary for the financial prediction component.

\subsection{Technical Stack}

\begin{itemize}
  \item \textbf{Frontend:} React 19, Vite, React Router, Zustand
  \item \textbf{Styling/UI:} Tailwind CSS and custom CSS
  \item \textbf{Backend:} Node.js and Express.js
  \item \textbf{Database:} MongoDB with Mongoose ODM
  \item \textbf{Authentication:} JWT, refresh tokens, bcryptjs
  \item \textbf{Validation:} Joi
  \item \textbf{Security middleware:} Helmet and express-rate-limit
  \item \textbf{File uploads:} Multer and Cloudinary
  \item \textbf{Frontend utilities:} Axios, React Hot Toast,
    Lucide React, QRCode React
  \item \textbf{ML service:} Python, Flask, Flask-CORS
  \item \textbf{Development/build tools:} Vite, npm and Oxlint
\end{itemize}

\section{Implementation Details}

\subsection{Authentication and Health ID Module}

The authentication module was implemented for both patients and doctors.
Patient registration stores the user's personal information and access
PIN, whereas doctor registration additionally stores specialization,
qualification, license number, experience, and optional hospital
association.

Passwords and patient access PINs are stored using bcrypt hashing.
Authentication uses JWT-based access tokens and refresh tokens. The
backend also applies rate limiting to authentication endpoints.

Each patient automatically receives a unique Health ID in the form

\[
\texttt{MB-YYYY-XXXXXX}
\]

during user creation. The Health ID is used throughout the platform to
identify patients and supports the doctor-access and emergency-access
workflows.

\subsection{Role-Based Access and Protected Routes}

The frontend contains dedicated route guards for authenticated users and
for role-based access. Patient and doctor pages are therefore separated
into protected workflows.

The backend similarly associates users with roles such as
\texttt{patient}, \texttt{doctor}, and \texttt{admin}. This provides a
foundation for limiting access to sensitive healthcare operations.

\subsection{Patient Profile and Emergency Profile}

The profile module supports retrieval and modification of patient
information including name, phone number, date of birth, gender, and
address. Profile image upload is supported through Cloudinary.

A separate emergency profile is maintained for information such as:

\begin{itemize}
  \item Blood group
  \item Allergies
  \item Chronic conditions
  \item Current medications
  \item Emergency contacts
  \item Insurance information
  \item Organ-donor status
  \item Additional emergency notes
\end{itemize}

The emergency profile is intentionally separated from the full medical
record so that an emergency-access workflow can expose a restricted
subset of information.

\subsection{Hospital and Doctor Discovery}

The healthcare discovery module implements hospital search with support
for:

\begin{itemize}
  \item Search by hospital name and related text fields
  \item City and state filtering
  \item Hospital type filtering
  \item Specialization filtering
  \item Emergency-availability filtering
  \item Minimum-rating filtering
  \item Pagination and sorting
\end{itemize}

Hospital records contain specializations, address information, ratings,
bed information, emergency availability, insurance information, and
verification state.

Hospital details additionally expose associated doctors. The doctor
model stores specialization, qualifications, experience, consultation
fee, availability, rating, and hospital association.

\subsection{Appointment Module}

The implementation extends the initial proposal with an appointment
workflow. Patients can submit appointments against a hospital by
providing:

\begin{itemize}
  \item Preferred date
  \item Specialist
  \item Preferred time period
\end{itemize}

Appointments initially enter a pending state. Doctor-side endpoints
support viewing appointments associated with the doctor's hospital and
updating appointment status to booked, rejected, or completed.

\subsection{Medical Records Module}

The medical-record subsystem provides centralized record storage for
patients.

Each record may contain:

\begin{itemize}
  \item Title and category
  \item Description and clinical notes
  \item Diagnosis
  \item Vital signs
  \item Hospital and doctor information
  \item Record date
  \item Tags
  \item Uploaded document metadata
\end{itemize}

Records may be uploaded by the patient, by a doctor, or imported through
the implemented mock hospital-fetch workflow.

Medical documents can be uploaded through Multer and stored using
Cloudinary. The backend provides record retrieval, filtering,
pagination, update, and deletion operations.

\subsection{Doctor Access and Consent Module}

A major implementation of the proposal's privacy objective is the
patient-controlled doctor-access mechanism.

A doctor first searches for the patient using the Health ID. The patient
then authorizes access using the four-digit access PIN.

On successful verification:

\begin{itemize}
  \item A doctor-patient access session is created.
  \item The session is valid for 24 hours.
  \item The access event is recorded in an audit log.
  \item The session contains permissions for records, prescriptions,
    emergency information, and creation of medical information.
\end{itemize}

Patients can inspect their active access sessions and revoke a session
before expiry. Access-log records therefore provide a history of access
grants and revocations.

The database also defines an expiry index on access sessions, allowing
temporary access to be represented as time-bounded rather than
permanent authorization.

\subsection{Prescription Module}

The prescription module stores prescriptions as structured records
rather than relying solely on uploaded documents.

A prescription can contain:

\begin{itemize}
  \item Diagnosis
  \item Symptoms
  \item Medicines
  \item Dosage and frequency
  \item Duration and timing
  \item Required tests
  \item Medical advice
  \item Follow-up date
  \item Relevant vitals
\end{itemize}

Doctors can create prescriptions for authorized patients, while patients
can retrieve their prescriptions and distinguish active prescriptions
from completed ones.

\subsection{Financial Assistance Module}

The financial-assistance subsystem implements the core proposal idea of
connecting healthcare costs with financial support.

A loan application stores treatment type, estimated treatment cost,
requested amount, income, existing EMIs, employment type, credit
history, tenure, and the resulting financial evaluation fields.

The implementation currently uses two complementary mechanisms.

\subsubsection{Rule-Based Eligibility}

A deterministic scoring function evaluates the application using factors
such as:

\begin{itemize}
  \item Requested amount relative to monthly income
  \item Debt-to-income condition
  \item Employment type
  \item Credit history
  \item A fixed age contribution in the current implementation
\end{itemize}

The rule score is accumulated using weighted criteria and is also used
for EMI calculation.

\subsubsection{ML Prediction Service}

A separate Flask service exposes a \texttt{/predict} API endpoint.
The backend sends application features to this service through Axios.

The current repository implementation uses the service as an MVP
prediction component. It computes a prediction score from income,
credit-history information, and loan amount and returns an approval
decision, score, maximum loan amount, interest-rate information, and a
reason.

The backend combines the ML response with the rule-based score to obtain
a demonstration-level final decision of approved, partial, or rejected.

This modular design allows the current prediction component to be
replaced by a trained model later without redesigning the main
healthcare backend.

\subsection{Fundraising Module}

The implementation also contains the proposed funding-gap workflow.

A patient can create a fundraising campaign containing:

\begin{itemize}
  \item Campaign title and story
  \item Total treatment requirement
  \item Loan-approved amount
  \item Funding gap
  \item Treatment and hospital information
  \item Supporting medical proof
  \item Campaign deadline
\end{itemize}

Donations are stored against campaigns and the amount raised and
remaining funding gap are updated accordingly.

The current implementation represents the payment operation as a
completed donation within the MVP rather than implementing a production
payment gateway. This leaves payment integration as a future
enhancement.

\subsection{Emergency and NFC-Oriented Workflow}

The application contains a public emergency route associated with a
patient's Health ID. The intended workflow is for an NFC card/tag to
encode an emergency-profile URL.

The architecture document specifies NFC write/read operations through
the Web NFC API and a public emergency page that exposes limited
information without login. The emergency information is intentionally
restricted and does not include complete medical records, prescriptions,
financial information, or the patient's full personal address.

\section{Current Implementation Status}

The repository at the mid-semester stage contains a functioning
full-stack application structure with significant portions of the
planned MVP implemented.

The implemented scope currently includes:

\begin{itemize}
  \item Patient and doctor registration
  \item Login and token-based authentication
  \item Automatic Health ID generation
  \item Patient profile management
  \item Emergency profile management
  \item Health ID page
  \item Hospital discovery and filtering
  \item Hospital detail pages
  \item Appointment creation and status handling
  \item Medical-record upload and management
  \item Doctor patient search
  \item Temporary PIN-based doctor access
  \item Access auditing and session revocation
  \item Structured prescriptions on the backend
  \item Loan application and eligibility workflow
  \item ML-service API integration
  \item Fundraiser and donation backend workflows
  \item Emergency public-access route
\end{itemize}

Some frontend routes are still represented by placeholder pages. In
particular, portions of the prescription, access-log, financial-result,
financial-history, fundraiser-management, and some doctor workflows
remain to be integrated fully into the UI.

% ==================================================
% CHAPTER 3
% ==================================================

\chapter{Results and Evaluation}

\section{Testing Methodology}

At the present stage, testing has focused primarily on verifying that
the implemented frontend, backend, database models, and service
interfaces work together correctly.

\subsection{Functional Testing}

The functional testing focus includes:

\begin{itemize}
  \item Patient registration followed by login
  \item Doctor registration followed by login
  \item Protected access to patient and doctor routes
  \item Patient profile retrieval and update
  \item Health ID generation and retrieval
  \item Hospital search and filtering
  \item Appointment creation and retrieval
  \item Medical-record creation and retrieval
  \item Doctor search for a patient by Health ID
  \item PIN verification and creation of a temporary access session
  \item Session revocation
  \item Prescription creation through doctor-side backend endpoints
  \item Financial eligibility evaluation
  \item Communication with the separate prediction API
  \item Fundraiser creation and donation updates
\end{itemize}

\subsection{Security-Oriented Testing}

Security-related implementation checks include:

\begin{itemize}
  \item Password hashing using bcrypt
  \item Access-PIN hashing using bcrypt
  \item JWT authentication
  \item Refresh-token validation
  \item Protected backend routes
  \item Role-based frontend route guards
  \item Rate limiting on authentication endpoints
  \item Helmet security middleware
  \item Validation through Joi schemas
  \item Restriction of patient-record queries to the authenticated
    patient or authorized doctor workflow
\end{itemize}

\subsection{Integration Testing}

Integration testing at this stage centers on communication between:

\begin{itemize}
  \item React frontend and Express REST API
  \item Express backend and MongoDB through Mongoose
  \item Express backend and Cloudinary
  \item Express backend and the Flask prediction service
\end{itemize}

\section{Performance Metrics}

The original project proposal defines evaluation criteria such as
patient workflow completion time, recommendation usefulness, document
access reliability, user satisfaction, and system availability.

At the mid-semester stage, the repository does not contain sufficient
measured experimental data to claim final numerical performance
targets. Therefore, this report records the implementation status
instead of inventing benchmark values.

\begin{table}[h]
\centering
\begin{tabular}{|p{0.38\textwidth}|p{0.20\textwidth}|p{0.28\textwidth}|}
  \hline
  \textbf{Evaluation Area} &
  \textbf{Target Stage} &
  \textbf{Current Status} \\
  \hline
  Authentication &
  Functional &
  Implemented for patient and doctor roles \\
  \hline
  Health ID &
  Functional &
  Automatic unique Health ID generation implemented \\
  \hline
  Hospital discovery &
  Functional &
  Search, filtering, sorting and pagination implemented \\
  \hline
  Medical records &
  Functional &
  CRUD, upload and mock hospital-fetch workflow implemented \\
  \hline
  Doctor access &
  Functional &
  PIN verification, 24-hour session and audit logs implemented \\
  \hline
  Appointments &
  Functional &
  Patient request and doctor status workflow implemented \\
  \hline
  Financial assistance &
  Prototype &
  Rule-based scoring and ML-service integration implemented \\
  \hline
  Emergency access &
  Prototype &
  Emergency profile and public Health-ID route implemented \\
  \hline
  Fundraising &
  Prototype &
  Campaign and donation backend workflow implemented \\
  \hline
\end{tabular}
\caption{Mid-Semester Implementation Status}
\label{tab:midsem-status}
\end{table}

\section{Analysis}

The project has progressed substantially beyond the initial conceptual
stage.

\subsection{Objectives Achieved}

The strongest progress has been made in the foundational healthcare
platform:

\begin{itemize}
  \item The proposed multi-role authentication system has been
    implemented.
  \item A centralized patient Health ID has been integrated into the
    data model and doctor-access workflow.
  \item Hospital and doctor discovery has moved from proposal stage to
    working backend and frontend modules.
  \item Medical records now have a persistent data model and file-upload
    workflow.
  \item The privacy objective has been implemented through temporary
    doctor access and access logging.
  \item Financial assistance is represented through both a rule-based
    engine and a separate prediction-service interface.
\end{itemize}

\subsection{Challenges Encountered}

The main challenges at the mid-semester stage are related to integration
and completion rather than the basic architecture.

First, the project contains multiple independent workflows that must
share authentication and patient identity securely. This increases the
need for consistent route protection and permission checks.

Second, healthcare information is sensitive. The design therefore has to
distinguish between general profile information, emergency information,
and complete medical records.

Third, the financial component combines deterministic rules with a
separate ML service. The backend must therefore continue functioning
even when the external model service is unavailable.

Finally, some backend functionality has been implemented ahead of the
corresponding frontend screens. This is visible in the repository where
several application routes currently use placeholder components.

\subsection{Current Limitations}

The following limitations are present in the current mid-semester
implementation:

\begin{itemize}
  \item Several frontend routes are still placeholders.
  \item The current Flask prediction endpoint is an MVP/demo prediction
    implementation rather than a fully trained production model.
  \item Production payment-gateway integration has not yet been
    implemented for donations.
  \item The project currently uses a mock hospital data-fetch workflow.
  \item Formal automated unit and integration test coverage is not yet
    represented in the repository at the level planned in the
    architecture specification.
  \item Final deployment, usability studies, and numerical performance
    evaluation remain future tasks.
\end{itemize}

% ==================================================
% CHAPTER 4
% ==================================================

\chapter{Conclusion}

\section{Summary of Work}

MediBridge was proposed as a centralized healthcare access and financial
assistance platform that would simplify healthcare discovery, medical
record management, financial-support workflows, and emergency access.

During the period leading to the mid-semester evaluation, the project
progressed from the proposal and architecture stage to a substantial
full-stack MVP implementation.

The developed system now contains a React frontend, Node.js/Express
REST backend, MongoDB persistence layer, cloud-based file-storage
integration, and a separate Flask prediction service.

The core healthcare workflow has been implemented through user
authentication, patient profiles, Health IDs, hospital discovery,
appointments, medical records, doctor interaction, privacy-controlled
access, and structured prescriptions. Financial assistance,
fundraising, and emergency-profile workflows have also been established
at the prototype level.

\section{Key Findings}

\begin{itemize}
  \item \textbf{The three-tier architecture is suitable for the project:}
    separating the React client, Express backend, and MongoDB persistence
    allows the different healthcare workflows to be developed
    modularly.

  \item \textbf{Patient identity is an effective integration point:}
    the MediBridge Health ID connects registration, doctor lookup,
    patient records, consent management, and the emergency-access
    workflow.

  \item \textbf{Temporary authorization improves privacy control:}
    PIN-based access sessions provide a concrete mechanism for granting
    doctors time-limited access rather than exposing patient records
    permanently.

  \item \textbf{Financial assistance benefits from modular design:}
    separating rule-based eligibility logic from the ML prediction API
    makes it possible to evolve the prediction component independently
    of the main platform.

  \item \textbf{The current implementation is ahead of the basic MVP
    foundation but not yet production-complete:} important work remains
    in frontend integration, automated testing, deployment, real-world
    integrations, and user evaluation.
\end{itemize}

\section{Future Work and Recommendations}

The next development stage will focus on completing the remaining
functional and integration work.

\begin{enumerate}
  \item Complete the frontend implementation of prescription viewing,
    access logs, active access sessions, loan results, loan history,
    fundraiser management, and remaining doctor workflows.

  \item Integrate the existing backend prescription and consent
    endpoints into their corresponding frontend pages.

  \item Replace the current demonstration prediction logic with the
    intended trained ML model and perform systematic model evaluation.

  \item Complete NFC write/read integration and validate the emergency
    access workflow on supported mobile devices.

  \item Add comprehensive unit and integration tests for authentication,
    access control, records, appointments, and financial workflows.

  \item Improve database indexing and query optimization as the
    healthcare dataset grows.

  \item Add production-grade donation/payment processing instead of the
    current MVP completion logic.

  \item Replace the mock hospital-fetch integration with suitable
    external healthcare interoperability mechanisms where feasible.

  \item Conduct usability testing with representative users and measure
    workflow completion time, document-access reliability,
    satisfaction, and system availability as defined in the original
    proposal.

  \item Complete cloud deployment and end-to-end validation of the
    platform.
\end{enumerate}

\section{Final Remarks}

The mid-semester implementation demonstrates that the proposed
MediBridge concept can be translated into a modular full-stack
software system. The foundational architecture and several of the most
important patient and doctor workflows are already implemented.

The remaining development effort is primarily focused on completing
frontend integration, strengthening testing and security validation,
refining the financial prediction component, and preparing the platform
for deployment and final evaluation.

% ==================================================
% BIBLIOGRAPHY
% ==================================================

\printbibliography[heading=bibintoc]

\end{document}